% ================= LỜI CẢM ƠN =================
\newpage
\thispagestyle{plain}

\begin{center}
    {\bfseries\Large LỜI CẢM ƠN}
\end{center}

\vspace{1cm}

Trước hết, tôi xin gửi lời cảm ơn chân thành đến quý thầy cô phụ trách học phần \coursename, đã
tận tình giảng dạy, định hướng và cung cấp những kiến thức nền tảng quan trọng cho đề tài này.

Những kiến thức về kiến trúc mạng nơ-ron tích chập, cơ chế chú ý và các phương pháp đánh giá mô
hình được truyền đạt trong học phần đã giúp tôi có cơ sở để tiếp cận đề tài
``\projecttitle''
một cách có hệ thống hơn, đặc biệt là ở khâu phân tích nguyên nhân đằng sau các con số thay vì chỉ
dừng lại ở việc báo cáo kết quả.

Tôi cũng xin cảm ơn các bạn cùng nhóm đã phối hợp trong quá trình xây dựng mã nguồn và tiến hành
thực nghiệm, cũng như đã trao đổi thẳng thắn về những kết quả không như kỳ vọng --- những trao đổi
đó đã góp phần định hình phần phân tích ở Chương~\ref{ch:ket_qua}.

Tôi xin cảm ơn \university{} và \faculty{} đã tạo điều kiện học tập, nghiên cứu và cung cấp môi
trường học thuật thuận lợi trong quá trình thực hiện bài tiểu luận này.

Do giới hạn về thời gian, tài nguyên tính toán và kinh nghiệm nghiên cứu, bài tiểu luận khó tránh
khỏi những thiếu sót; một số hạn chế đã được tôi nêu rõ ở Chương~\ref{ch:tong_ket}. Tôi rất mong
nhận được những góp ý từ giảng viên để có thể tiếp tục hoàn thiện nội dung và phát triển đề tài
tốt hơn trong các nghiên cứu tiếp theo.
